\documentclass[11pt,a4paper]{article}
\usepackage[margin=27mm]{geometry}
\usepackage{amsmath,amssymb,amsthm,booktabs,hyperref}
\newtheorem{theorem}{Theorem}
\newtheorem{lemma}[theorem]{Lemma}
\newtheorem{corollary}[theorem]{Corollary}
\newcommand{\F}{\mathbb F}
\newcommand{\Sz}{\operatorname{Sz}}
\newcommand{\gen}[1]{\langle #1\rangle}
\newcommand{\N}{\operatorname{N}}
\newcommand{\Tr}{\operatorname{Tr}}
\emergencystretch=1em
\title{Subfield parabolics and the corrected Suzuki classification}
\author{Human Compression project}\date{}
\begin{document}\maketitle
\begin{abstract}
We complete the subgroup argument for configurations with a maximal
subfield parabolic. A normalizer condition in a finite two-group forces
every cover of a subfield point stabilizer into the central overgroup
already treated by the correspondence with scalar subspaces. Thus that
correspondence supplies the entire corrected family. We count its field
orbits by an explicit formula for fixed projective lines under Frobenius.
Together with the separately established normalizer families, this gives
a corrected classification and a finite arithmetic formula for the total
number of configurations. The numerical source tables and diagrams still
require their separate complete audit; this draft does not by itself
claim whole-source completion.
\end{abstract}

\section{Acquired results and notation}
Let $q=2^n$, where $n\ge3$ is odd, and let $G=\Sz(q)$. Write
$a_\pm(q)=q\pm2^{(n+1)/2}+1$. We use the classical Suzuki
partition, subgroup classification, ovoid action and automorphism group
specified in the accepted note \emph{Normalizer fibres in rank-two
Suzuki geometries}. In particular every proper nonsolvable subgroup is
a simple subfield group, all subfield groups of a prescribed order are
conjugate, and their proper maximal simple subfield groups have prime
relative field degree. The solvable maximal subgroups have already
been treated, and a point stabilizer cannot be a parabolic in the
geometries under consideration.

Fix a prime $r$ dividing $n$, put $d=n/r\ge3$, $s=2^d$,
$F=\F_q$, $K=\F_s$, and let $S=\Sz(s)$ be the standard maximal
subfield subgroup. Its point stabilizer is
$B_s=U_KH_K$, with $U_K=\{u(a,b):a,b\in K\}$ and
$H_K=\{h(k):k\in K^*\}$. We retain the same faithful root-group
coordinates as in the preceding notes:
\[
 u(a,b)u(c,e)=u(a+c,b+e+ac^\theta),\quad
 u(a,b)^{h(k)}=u(k^\theta a,k^{\theta+2}b),
 \qquad\theta=2^{(n+1)/2}.
                                                        \tag{1}
\]
Put $P=U(F,F)$ and $Z_F=U(0,F)$. The two coordinate actions of
$K^*$ are the full scalar group of $K$. The automorphism
$x\mapsto x^\theta$ on $K$ has fixed field $\F_2$.

The accepted scalar-subspace note proves that the overgroups of
$B_s$ in $Z_FB_s$ are exactly
\[
 Q_T=Z_TB_s,\qquad K\le T\le F\text{ a }K\text{-subspace},
 \quad Z_T=\{u(0,t):t\in T\}.                            \tag{2}
\]
A group $Q_T$ covers $B_s$ precisely when $\dim_KT=2$; it then
has order $s^3(s-1)$. The geometries $\Gamma(G;S,Q_T)$ are firm,
connected, flag transitive, residually weakly primitive and locally
doubly transitive. Their equivalence classes, in the source's convention
of automorphisms of the specified group with type interchange allowed,
are exactly the Frobenius orbits on the projective space
$\mathbb P_K(F/K)$. These statements are reused, not proved twice.
The task here is to prove exhaustion and compute those orbits.

\section{Every point-stabilizer cover is central}
\begin{lemma}\label{normalizer}
Inside the point stabilizer $B_q=P\rtimes F^*$,
\[
 \N_P(U_K)=Z_FU_K,\qquad \N_{B_q}(U_K)=Z_FU_KH_K.
                                                        \tag{3}
\]
\end{lemma}
\begin{proof}
An element $u(a,b)\in P$ normalizes $U_K$ precisely when
$ac^\theta+ca^\theta\in K$ for every $c\in K$, by (1).
Taking $c=1$ gives $a+a^\theta\in K$. Subtracting $c$ times
this relation from the general one gives
$a(c^\theta+c)\in K$. Choose $c\in K\setminus\F_2$; its
coefficient is nonzero, so $a\in K$. There is no restriction
on $b$. This proves the first equality.

Conjugation by $P$ does not change first coordinates. Therefore,
if $u(a,b)h(t)$ normalizes $U_K$, its action on the first-coordinate
image $K$ gives $t^\theta K=K$, hence $t\in K^*$.
Removing this diagonal factor reduces to the first equality.
Conversely, all displayed elements normalize $U_K$.
\end{proof}

\begin{theorem}\label{covers}
If $B_s<L<B_q$ and $B_s$ is maximal in $L$, then $L=Q_T$ for
a two-dimensional $K$-subspace $T$ containing $K$. These are all
proper subgroups outside $S$ in which $B_s$ is maximal.
\end{theorem}
\begin{proof}
Put $V=L\cap P$, a normal two-subgroup of $L$ containing $U_K$.
If $V=U_K$, then $L\le\N_{B_q}(U_K)$. Equation (3) makes its
image in $B_q/P$ equal to $H_K$. Because $H_K\le L$, we obtain
$L=VH_K=B_s$, a contradiction. Thus $V>U_K$.

The normalizer condition in a finite $p$-group gives
$\N_V(U_K)>U_K$. For completeness, a proper subgroup $H$ of
a finite $p$-group acts on its cosets; the fixed cosets are
$\N(H)/H$, all other orbits have length divisible by $p$, and
the total index is divisible by $p$. There is a fixed coset, so
there are at least $p$ of them. This proves the condition.
Now $H_K$ normalizes both $V$ and $U_K$, hence normalizes
$\N_V(U_K)$. Therefore
\[
             B_s<\N_V(U_K)H_K\le L.
\]
Maximality forces equality with $L$. Consequently $U_K$ is
normal in $L$, and (3) gives $L\le Z_FB_s$. The acquired
correspondence (2) proves the claimed form and dimension.

Finally consider any proper overgroup $L$ of $B_s$ outside $S$.
By the classical subgroup classification it either lies in a
solvable normalizer or is a simple subfield group. Its large
two-subgroup $U_K$ excludes the cyclic torus normalizers.
If it lies in a point stabilizer, a nonidentity involution of
$U_K$ forces that stabilizer to be the particular $B_q$.
If $L=\Sz(t)$ and $B_s$ is maximal in $L$, the classical maximal
list forces $B_s$ to be a point stabilizer of $L$, hence $t=s$.
The accepted uniqueness theorem for a subfield overgroup of
$B_s$ would give $L=S$. Thus this case is excluded and the
list already obtained is exhaustive.
\end{proof}

This strengthens the accepted correction: the groups $Q_T$ are
the entire replacement for the first row of source Table 8, not
only a family contradicting that row. The printed group of order
$qs(s-1)$ does not cover $B_s$ when $r\ge3$; the correct groups
have order $s^3(s-1)$ and are parametrized as in (2).

\section{The other subfield intersections}
We shall use a further consequence of the normalizer incidence
formula. If $\F_t\subseteq\F_s\subseteq\F_q$, every fixed
$\Sz(t)$ is in exactly one $\Sz(s)$ of the
prescribed order. Here and below both smaller field orders are at
least eight. Indeed, all the subfield groups of each order are
conjugate, both are self-normalizing by simplicity and the classical
subgroup classification, and the subgroups of order $|\Sz(t)|$
inside $\Sz(s)$ form one conjugacy class. The incidence formula
therefore gives one overgroup. This is a uniqueness assertion for
each prescribed field order, not a claim that the entire interval
has only one intermediate subgroup.

\begin{theorem}\label{list}
Besides the central covers in Theorem~\ref{covers}, all configurations
with maximal parabolic $S=\Sz(s)$ are as follows:
\[
\begin{array}{c|c|c}
 J=S\cap L&L&\text{condition}\\\hline
 D_{2(s-1)}&D_{2m}&m/(s-1)\text{ prime},\ m\mid q-1\\
 C_{a_\delta(s)}\rtimes C_4&C_m\rtimes C_4&
 m/a_\delta(s)\text{ prime},\ m\mid a_\epsilon(q)\\
 \Sz(2^e)&\Sz(2^{ep})&
 n=erp,\ r,p\text{ distinct primes},\ e\ge3.
\end{array}                                               \tag{4}
\]
In the second row $\delta\in\{+1,-1\}$ and the ambient sign
$\epsilon$ is the unique one for which
$a_\delta(s)\mid a_\epsilon(q)$. Each permitted numerical choice
in the first two rows gives one equivalence class. Each unordered
pair of distinct prime indices in the last row gives one class.
\end{theorem}
\begin{proof}
The maximal-subgroup list in $S$ leaves precisely the three
intersection types in (4), after its point stabilizer.
For $J=D_{2(s-1)}$, its odd rotation subgroup determines the
unique split torus of $G$. An overgroup in a point stabilizer
cannot contain the involution inverting that nontrivial odd
subgroup: projection to the point stabilizer's abelian odd
quotient would make inversion trivial. Coprimality excludes
the nonsplit normalizers. Thus a solvable overgroup is in the
unique split dihedral normalizer, and the cyclic maximal-subgroup
calculation gives exactly the first row. A simple subfield
overgroup in which $J$ is maximal must have field order $s$;
the already accepted uniqueness theorem would make it $S$.
Within the ambient torus normalizer the group of each displayed
order containing this fixed $J$ is unique, so each choice gives
one configuration class.

For $J=C_{a_\delta(s)}\rtimes C_4$, its characteristic odd
subgroup similarly determines its unique ambient nonsplit
torus. The point stabilizer and split normalizer are excluded
by their odd orders. Its solvable covers are therefore exactly
the prime torus extensions in the second row. A simple overgroup
in which $J$ is maximal must again have field order $s$:
the two functions $a_+(2^d)$ and $a_-(2^d)$ strictly increase
through nonoverlapping intervals as the odd exponent $d\ge3$
increases, and the two signs at a fixed exponent are distinct.
Hence its odd torus order determines $s$ and $\delta$.
The uniqueness theorem then excludes it. For a fixed $J$ each
allowed $m$ gives one group, proving the count one.

If $J=\Sz(2^e)$ is a simple subfield group, then $d/e=p$
is prime, by maximality in $S$. An overgroup $L$ is nonsolvable,
so is a subfield group $\Sz(2^{e\ell})$, with $\ell$ prime
by maximality of $J$ in $L$. Since $n/e=rp$, necessarily
$\ell=r$ or $p$. If $\ell=p$, uniqueness gives $L=S$.
Thus $\ell=r\ne p$. Conversely distinct $r,p$ give two
different maximal subfield groups over the fixed $J$. Their
intersection is exactly $J$: it contains $J$, maximal in each,
and neither contains the other. Their conjugacy class is unique
by the same incidence argument, and type interchange identifies
the two orders of the prime pair.

For the first two rows, the proposed $L$ lies outside $S$ and
$J$ is maximal in both; hence the intersection is exactly $J$.
In every row the maximality of $S$ gives generation of $G$ and
the rank-two coset criterion supplies firmness, connectivity,
flag transitivity and residual weak primitivity.
\end{proof}

The nonoverlap used in the proof is elementary. For $s\ge8$,
$a_+(s)\le3s/2+1$, whereas
$a_-(4s)=4s-2\sqrt{2s}+1>3s/2+1$. Adjacent odd exponents
multiply the field order by four, proving the assertion.
The case $2^e=2$ is not included in the simple-subfield row:
$\Sz(2)$ is the solvable group $C_5\rtimes C_4$, and any such
maximal intersection is already governed by the nonsplit row.

\section{Frobenius orbits of the central covers}
The following argument is finite field algebra; it is independent
of the subgroup classification. It applies to every odd $d\ge1$
and odd prime $r$, with $n=dr$, even when $r$ divides $d$.

\begin{lemma}\label{hilbert}
For a cyclic finite-field extension $E/D$ with generator $\tau$,
the elements of norm one are exactly $\tau(c)/c$ for $c\in E^*$.
If the extension degree is odd in characteristic two, invariants
are exact on every sequence of vector spaces with compatible
actions: an invariant vector in a quotient has an invariant lift.
\end{lemma}
\begin{proof}
The map $c\mapsto\tau(c)/c$ has kernel $D^*$, and its image
has size $(|E|-1)/(|D|-1)$. The norm map on the cyclic group
$E^*$ is the power $(|E|-1)/(|D|-1)$, is onto $D^*$, and
has a kernel of that same size. The image lies in the kernel
by telescoping, proving equality. For the second assertion,
average any lift over the acting cyclic group. Its odd order
is invertible in characteristic two, and averaging preserves
the specified quotient vector.
\end{proof}

\begin{theorem}\label{burnside}
For $0\le j<n$, put $h=\gcd(d,j)$ and $g=\gcd(n,j)$,
with the usual convention $\gcd(a,0)=a$. The number $b_j$ of
$K$-projective lines in $F/K$ fixed by $\phi^j$ is
\[
 b_j=\begin{cases}
 \displaystyle\frac{2^{h(r-1)}-1}{2^h-1},&g=rh,\\[3mm]
 \gcd(r,2^h-1)-1,&g=h.
 \end{cases}
                                                        \tag{5}
\]
These are the only cases. Consequently the number of central
cover classes is
\[
                         C(n,d)=\frac1n\sum_{j=0}^{n-1}b_j.
                                                        \tag{6}
\]
\end{theorem}
\begin{proof}
Set $\sigma=\phi^j$ and $m=d/h$, the order of its restriction
to $K$. Since $\gcd(d/h,j/h)=1$,
\[
       g=h\gcd(r,j/h)\in\{h,rh\}.
\]
This identity does not require $r$ to be coprime to $d$.
The order of $\sigma$ on $F$ is $n/g$.

Suppose $g=rh$. Then $\sigma$ has order $m$ on both fields.
The invariant fields are $F^\sigma=\F_{2^{rh}}$ and
$K^\sigma=\F_{2^h}$. Exactness by averaging gives
\[
              |(F/K)^\sigma|=2^{h(r-1)}.
\]
On a fixed $K$-line choose $v\ne0$ with $\sigma(v)=\lambda v$.
Applying $\sigma^m=1$ gives $N_{K/\F_{2^h}}(\lambda)=1$.
Lemma~\ref{hilbert} permits a scalar change making $v$ fixed.
Each fixed line then has exactly $2^h-1$ nonzero fixed vectors.
Dividing gives the first expression in (5).

Suppose $g=h$. The automorphism
$\rho=\sigma^m=\phi^{jd/h}$ has order $r$ on $F$ and fixes
$K$. Its one-eigenspace in $F$ is exactly $K$.
Since $r$ is odd, trace identifies $F/K$ with the trace-zero
complement $V=\ker(\Tr_{F/K})$; this complement is also
$\sigma$-invariant. The one-eigenspace is removed, so $V^\rho=0$.
On a $\sigma$-fixed $K$-line, $\sigma$ acts semilinearly and
$\rho$ acts as a scalar
$\mu=N_{K/\F_{2^h}}(\lambda)\in\F_{2^h}^*$.
Necessarily $\mu^r=1$ and $\mu\ne1$.

Conversely, for every such $\mu$, Lemma~\ref{hilbert} applied
to $F/K$ gives a nonzero solution of $\rho(v)=\mu v$.
Its solutions form one $K$-line: the ratio of two nonzero
solutions is fixed by $\rho$. For $\mu\ne1$ this line has
trace zero, so lies in $V$. Commutation of $\sigma$ and $\rho$
sends the $\mu$-eigenline to the $\sigma(\mu)$-eigenline.
Thus the line is $\sigma$-stable exactly when
$\mu\in\F_{2^h}$. There are $\gcd(r,2^h-1)-1$ such nontrivial
roots. This proves the second expression. Finally, counting
the pairs consisting of an orbit and an automorphism fixing
one of its points gives Burnside's formula (6).
\end{proof}

For example,
\[
 C(9,3)=1,\quad C(15,3)=41,\quad C(15,5)=3,\quad
 C(21,3)=1791,\quad C(21,7)=7.
\]
The first two agree with the independently accepted literal
projective-orbit certificates in the scalar-subspace note.
The formula also covers the repeated prime index cases, such
as $C(27,9)=19$; no coprimality of the absolute and relative
field degrees has been imposed.

\section{A corrected global list and total}
Let $\omega(a)$ denote the number of distinct prime divisors of
the positive integer $a$, with $\omega(1)=0$. Each simple
subfield case above has maximal absolute field exponent
$d=n/r\ge3$, for a prime divisor $r$ of $n$. For each sign
$\delta$, let $t_\delta(q,s)$ be the unique ambient nonsplit
torus order divisible by $a_\delta(s)$.

\begin{theorem}\label{global}
The complete classification consists of the solvable maximal
parabolic families in the accepted normalizer note, together
with (2) and (4), identifying a configuration when both descriptions
give the same unordered pair. If $n$ is composite, its total is
\begin{equation}
 T(n)=1+E(n)+\sum_{\substack{r\mid n\ \text{prime}\\d=n/r\ge3}}
                         [C(n,d)+A(n,d)],               \tag{7}
\end{equation}
where
\begin{align*}
 A(n,d)&=\omega\!\left(\frac{q-1}{2^d-1}\right)
 +\sum_{\delta=\pm1}\omega\!\left(
     \frac{t_\delta(q,2^d)}{a_\delta(2^d)}\right),\\
 E(n)&=\#\{\{r,p\}:r\ne p\text{ prime},\ rp\mid n,\ n/(rp)\ge3\}.
\end{align*}
If $n$ is prime, put
\[
 v=\frac{q-2}{2n},\qquad
 I=1+\frac{(q-2)(q+4)}{4n},\qquad
 H=1+\frac{(q-2)(q+2)}{4n}.
\]
Let $\eta_0,\eta_+,\eta_-$ be the indicators that respectively
$q-1,a_+(q),a_-(q)$ are prime. Then
\begin{align}
 T(n)={}&1+\eta_0\bigl[I+2v+1+
                 H(\omega(a_+(q))+\omega(a_-(q)))\bigr]\notag\\
 &+\sum_{\epsilon=\pm1}\eta_\epsilon
       \bigl[2v+(v+1)\omega(a_{-\epsilon}(q))\bigr]
       -\eta_+\eta_-(v+1).                              \tag{8}
\end{align}
\end{theorem}
\begin{proof}
Residual weak primitivity gives at least one maximal parabolic
in $G$. The accepted preliminary obstruction excludes a point
stabilizer. The accepted normalizer classification treats the
three torus normalizers; Theorems~\ref{covers} and \ref{list}
treat every maximal intersection in a subfield parabolic.
These exhaust the classical maximal-subgroup list, proving the
classification.

If $n$ is composite, all three ambient odd factors are composite.
The universal pair $(D_{2(q-1)},q\!:\!(q-1))$ contributes one.
All other families have a subfield parabolic. For a fixed $s$,
the first two rows of (4) are counted by the distinct prime
divisors of the indicated quotients, and the central family by
$C(n,d)$. The last row has two maximal subfield parabolics and
must be counted once for each unordered pair of distinct prime
indices. Their simple intersection must have field order at
least eight. Distinct field orders or distinct torus orders do
not identify these families. Including the full ambient torus
normalizer among the possible $L$ in (4) already includes the
corresponding family from the normalizer note, so it is not
added again. This proves (7).

If $n$ is prime there are no proper simple subfield groups.
The split prime case adds the same-type contribution $I$,
the four-group and cyclic-four contributions $v$ and $v+1$,
and one $H$ for each prime divisor of either nonsplit order.
For each prime nonsplit order there are two contributions $v$
and one $v+1$ per prime divisor of the opposite order. When
both nonsplit orders are prime, their common mixed family has
been counted twice; subtract $v+1$. No other families overlap.
This is (8).
\end{proof}

\section{Local double transitivity and source disposition}
The universal pair is locally doubly transitive: its local degrees
are two and $q$, and the latter action is the affine action of
$\F_q\rtimes\F_q^*$. Every corrected central cover is also
locally doubly transitive by the acquired scalar-subspace theorem.
Thus the source's claim that only the universal family survives
this condition for $n>5$ requires correction when $n$ is composite.

There are no further locally doubly transitive families except,
at $q=8$, the pairs $(C_5\!:\!C_4,C_5\!:\!C_4)$ and
$(C_5\!:\!C_4,C_4\times C_2)$. Here a cyclic complement of
order four acts regularly on the four nonzero elements of $C_5$,
and the other local action in the second pair has degree two.
For the prime split families, an intersection of order two
cannot be transitive on the $q-2\ge6$ other cosets. For a prime
nonsplit normalizer with intersection $C_4$, double transitivity
requires its prime torus order to satisfy $p-1\le4$, so $p=5$;
the opposite torus cannot also have order five. For the remaining
subfield families the local action of $S$ already fails double
transitivity. If $J$ is a torus normalizer of $S$, then
$|J|\le8s$ and $|S|>|J|(|J|+1)$ for $s\ge8$, so
$[S:J]-1>|J|$. If $J=\Sz(t)$ is a proper simple subfield group,
then $s\ge t^3$, $|J|<t^5$ and $|S|>s^5/2>t^{10}+t^5$;
the same inequality applies. A point stabilizer of size $|J|$
cannot act transitively on more than $|J|$ points.

This supplies a corrected form of source Lemma 6.9 and its
Table 8, the complete family statement of Theorem 7.1, and
the corresponding local double-transitivity column in Tables
9--11. The source's numerical Table 12 and its 29 small-group
diagrams need explicit numerical and invariant joins. They
are not silently certified by a general formula. The original
source is D.~Leemans, \emph{The rank 2 geometries of the simple
Suzuki groups $\Sz(q)$}, Beitr\"age zur Algebra und Geometrie
39 (1998), no.~1, 97--120. Its retained structural inputs remain
attributed to Suzuki, Tits and L\"uneburg. The earlier component's
phrase ``right coset $gJ$'' should read simply ``coset $gJ$'';
its equations are unaffected. No Lean or whole-original completion
claim is made by this draft.
\end{document}
